\chapter{The Synonym Mode}\index{modes!Synonym Mode|textbf}
\label{ch:synonyms}

\begin{versebox}
Many are the synonyms ---\\
one thing with a hundred names.\\
Know the names, and know the thing:\\
the teaching stays the same.
\end{versebox}

\noindent The Buddha had a remarkable habit. He would describe the same thing over and over, using different words each time. Not because he was being repetitive, but because each synonym illuminated a different facet of the same reality --- like turning a diamond in the light.\footnote{Pali: \textit{``Tattha katamo vevacano hāro? `Vevacanāni bahūnī'ti. Yathā ekaṃ bhagavā dhammaṃ aññamaññehi vevacanehi niddisati.''} The mnemonic fragment is ``many are the synonyms'' (\textit{vevacanāni bahūni}). The word \textit{vevacana} means ``synonym, alternative expression.'' The Synonym Mode teaches that when the Buddha uses multiple terms, they often point to a single reality --- and knowing this prevents the reader from multiplying entities unnecessarily.}

The Synonym Mode is the tool for recognizing when different words mean the same thing\index{synonyms|textbf}. Without it, you might read a sutta that mentions ``hope,'' ``longing,'' ``delight,'' and ``the current'' and think the Buddha was talking about four different mental states. He was talking about one: craving\index{craving|textbf}.

\section*{A Hundred Names for Craving}

The demonstration begins with a verse we have already met:

\begin{versebox}
Hope, longing, and delight ---\\
currents lodged in many realms ---\\
born from the root of ignorance,\\
all yearnings I have ended, root and stem.
\end{versebox}

The Netti now takes this verse apart word by word.\footnote{Pali: \textit{``Āsā ca pihā abhinandanā ca, anekadhātūsu sarā patiṭṭhitā; / Aññāṇamūlappabhavā pajappitā, sabbā mayā byantikatā samūlikā''ti.} (Cf.\ SN IV.409.) The same verse used in Chapter~\ref{ch:fitness} to illustrate the fire-analogy. Here the Netti uses it differently: not to show the unity of craving's characteristic (clinging), but to map its synonyms. \textbf{A note on our approach}: the original Pali provides detailed definitions for each term in the verse, then catalogues the eighteen elements\index{elements (eighteen)|textbf}, then maps the ``currents'' onto the six sense-objects. We have woven this material into a flowing narrative, with the technical details in the notes.}

``Hope'' is the longing for something in the future --- the feeling that what you want \textit{will} come. ``Longing'' is desire for something happening right now --- you see someone who has what you want, and you think: \textit{I wish I were like that}. ``Delight'' is the savoring of what has already arrived --- clinging to the pleasant, rejoicing in the agreeable.\footnote{Pali: \textit{``Āsā nāma vuccati yā bhavissassa atthassa āsīsanā avassaṃ āgamissatīti āsāssa uppajjati. Pihā nāma yā vattamānassa atthassa patthanā, seyyataraṃ vā disvā `ediso bhaveyya'nti pihāssa uppajjati. Atthanipphattipaṭipālanā abhinandanā nāma, piyaṃ vā ñātiṃ abhinandati, piyaṃ vā dhammaṃ abhinandati, appaṭikūlato vā abhinandati.''} Three synonyms for craving, three temporal orientations: \textit{āsā} (hope) = future-directed; \textit{pihā} (longing) = present-directed; \textit{abhinandanā} (delight) = oriented toward what has been attained. Same mental state, different time-horizons.}

Three words. Three temporal angles on the same thing. Hope looks forward. Longing looks sideways at what others have. Delight looks at what you're holding. But all three are craving.

``The many realms'' are the eighteen elements of experience --- the eye and visible form and eye-consciousness, the ear and sound and ear-consciousness, and so on through all six senses. These are the realms in which the ``currents'' flow.\footnote{Pali: \textit{``Anekadhātūti cakkhudhātu rūpadhātu cakkhuviññāṇadhātu, sotadhātu saddadhātu sotaviññāṇadhātu \ldots manodhātu dhammadhātu manoviññāṇadhātu.''} The eighteen elements (\textit{aṭṭhārasa dhātu}): six sense-organs $\times$ three (organ + object + consciousness) = eighteen. These are the ``many realms'' in which craving takes up residence.}

And ``the currents''? These are the tendencies --- some people are drawn to sights, some to sounds, some to smells, some to tastes, some to touch, some to ideas. The currents flow wherever there is a sense-object to flow toward. And when you add up all the feelings that arise in these encounters --- the six types of household grief, the six types of household joy, the six types of renunciation-grief, and the six types of renunciation-joy --- you get twenty-four terms. All of them synonyms for craving.\footnote{Pali: \textit{``Sarā'ti keci rūpādhimuttā keci saddādhimuttā \ldots Tattha yāni cha gehasitāni domanassāni yāni ca cha gehasitāni somanassāni yāni ca cha nekkhammasitāni domanassāni yāni ca cha nekkhammasitāni somanassāni, imāni catuvīsapadāni taṇhāpakkho, taṇhāya etaṃ vevacanaṃ.''} Twenty-four synonyms for craving, derived from the Saḷāyatana Vibhaṅga Sutta (MN 137): 6 household griefs + 6 household joys + 6 renunciation-griefs + 6 renunciation-joys = 24. The six household-based equanimities are classified separately as belonging to the ``view-side'' (\textit{diṭṭhipakkho}).}

And there are still more. ``Delight in the teaching, affection for the teaching, clinging to the teaching'' --- even these, which sound wholesome, are synonyms for craving when they involve attachment.\footnote{Pali: \textit{``Sāyeva patthanākārena dhammanandī dhammapemaṃ dhammajjhosānan'ti taṇhāya etaṃ vevacanaṃ.''} A crucial point: even ``delight in the Dhamma'' (\textit{dhammanandī}) can be craving if it involves clinging. The Synonym Mode helps you recognize craving even in its most sophisticated disguises.}

\section*{Three Names for the Mind}

The mind gets three names too: mind as consciousness, mind as thought, mind as cognition. Same mental event, three angles.\footnote{Pali: \textit{``Cittaṃ mano viññāṇan'ti cittassa etaṃ vevacanaṃ.''} The three standard terms for mind: \textit{citta} (mind as heart, as subjective experience), \textit{mano} (mind as thinking, as intellect), \textit{viññāṇa} (mind as cognizing, as awareness). All three refer to the same reality.} And the mind-faculty, the mind-element, the mind-base, and ``knowing'' --- all synonyms for the mind.\footnote{Pali: \textit{``Manindriyaṃ manodhātu manāyatanaṃ vijānanā'ti manassetaṃ vevacanaṃ.''}}

\section*{A Constellation of Wisdom}

But the most impressive list of synonyms belongs to wisdom. The Netti gives at least twenty: the wisdom-faculty, the wisdom-power, higher wisdom, the training in wisdom, the wisdom-aggregate, the awakening-factor of investigation, knowledge, right view, discernment, insight, knowledge of the teaching, knowledge of meaning, inferential knowledge, knowledge of destruction, knowledge of non-arising, the faculty of ``I shall know what is not yet known,'' the faculty of ``knowing,'' the faculty of ``one who has known,'' the eye, true knowledge, intelligence, understanding, sagacity, light.\footnote{Pali: \textit{``Paññindriyaṃ paññābalaṃ adhipaññā sikkhā paññā paññākkhandho dhammavicayasambojjhaṅgo ñāṇaṃ sammādiṭṭhi tīraṇā vipassanā dhamme ñāṇaṃ atthe ñāṇaṃ anvaye ñāṇaṃ khaye ñāṇaṃ anuppāde ñāṇaṃ anaññātaññassāmītindriyaṃ aññindriyaṃ aññātāvindriyaṃ cakkhu vijjā buddhi bhūri medhā āloko, yaṃ vā pana yaṃ kiñci aññaṃpi evaṃ jātiyaṃ, paññāya etaṃ vevacanaṃ.''} At least twenty-four synonyms for wisdom (\textit{paññā}). The concluding phrase --- \textit{``yaṃ vā pana yaṃ kiñci aññaṃpi evaṃ jātiyaṃ''} (``and whatever else of the same kind'') --- signals that even this long list is not exhaustive.}

All of them: wisdom. One thing, seen from two dozen angles. The Netti adds: all five supramundane faculties are, at their core, wisdom --- but each has its own emphasis. Faith is wisdom in its aspect of commitment. Energy is wisdom in its aspect of effort. Mindfulness is wisdom in its aspect of non-floating. Concentration is wisdom in its aspect of non-distraction. And wisdom itself is wisdom in its aspect of understanding.\footnote{Pali: \textit{``Pañcindriyāni lokuttarāni, sabbā paññā. Api ca ādhipateyyaṭṭhena saddhā, ārambhaṭṭhena vīriyaṃ, apilāpanaṭṭhena sati, avikkhepaṭṭhena samādhi, pajānanaṭṭhena paññā.''} A remarkable claim: the five supramundane faculties are all ``wisdom'' (\textit{sabbā paññā}), differentiated only by emphasis. Faith is wisdom-as-commitment; energy is wisdom-as-effort; mindfulness is wisdom-as-staying-present; concentration is wisdom-as-non-distraction; wisdom is wisdom-as-understanding.}

\section*{Synonyms for the Buddha, the Teaching, and the Community}

The Synonym Mode then catalogues the many names used in the six recollections --- the traditional Buddhist meditation on the qualities of the Buddha, the Teaching, the Community, virtue, generosity, and the gods.

The Buddha: accomplished, fully self-awakened, perfect in knowledge and conduct, well-gone, knower of the worlds, unsurpassed trainer of persons to be tamed, teacher of gods and humans, awakened, blessed. But also: one who has reached the power of the ten powers, attained the four confidences, achieved the four analytical knowledges, abandoned the four yokes, gone beyond wrong courses, pulled out the dart, healed the wound, crushed the thorn, extinguished the obsessions, gone beyond bondage, untangled the knots, passed beyond the underlying tendencies, broken the darkness, endowed with vision, gone beyond worldly states.\footnote{Pali: \textit{``Yathā ca buddhānussatiyaṃ vuttaṃ itipi so bhagavā arahaṃ sammāsambuddho vijjācaraṇasampanno sugato lokavidū anuttaro purisadammasārathi satthā devamanussānaṃ buddho bhagavā. Balanipphattigato vesārajjappatto adhigatappaṭisambhido \ldots ābhaṃkaro pabhaṃkaro dhammobhāsapajjotakaroti.''} Two layers of synonyms: the standard nine epithets of the Buddha (from the canonical \textit{buddhānussati} formula), followed by an extended list of additional designations. Some are martial metaphors (crushed the thorn, pulled out the dart), others visual (broken the darkness, torch-bearer, light-maker, dawn-bringer). All point to one person.}

The Teaching: well-proclaimed, visible here and now, immediately effective, inviting investigation, leading onward, to be personally known by the wise. But also: the crushing of intoxication, the quenching of thirst, the uprooting of attachment, the cutting of the round, the empty, the supremely rare, the destruction of craving, dispassion, cessation, nirvana.\footnote{Pali: \textit{``Svākkhāto bhagavatā dhammo sandiṭṭhiko akāliko ehipassiko opaneyyiko paccattaṃ veditabbo viññūhi. Yadidaṃ madanimmadano pipāsavinayo ālayasamugghāṭo vaṭṭūpacchedo suññato atidullabho taṇhakkhayo virāgo nirodho nibbānaṃ.''} The standard dhammānussati formula, extended. Note the progression from practical qualities (visible, inviting investigation) to ever-deeper synonyms for nirvana\index{nibbāna|textbf} itself.}

And the Netti caps this with an extraordinary verse --- a cascade of synonyms for nirvana, the unconditioned:\footnote{Pali verses from SN IV.409 and related sources. The Netti presents five stanzas, each piling synonym upon synonym for \textit{nibbāna}. We translate the core stanzas here.}

\begin{versebox}
Unconditioned, boundless, uncorrupted ---\\
truth, the far shore, subtle, so hard to see.\\
Unaging, stable, without appearance,\\
non-proliferating --- peace.

The deathless, the exquisite, the secure, the safe ---\\
destruction of craving, marvelous, astonishing.\\
Free from calamity, calamity's cessation ---\\
this is nirvana, taught by the Fortunate One.

Unborn, unbecome, untroubled,\\
uncreated, sorrowless, beyond sorrow.\\
Free from harm, harm's cessation ---\\
this is nirvana, taught by the Fortunate One.

Deep and hard to see,\\
surpassing yet unsurpassed ---\\
peerless, unequalled,\\
the highest of the highest.

A shelter, a refuge, a retreat, the stainless ---\\
flawless and pure, they say.\\
An island, happiness, the immeasurable, a foundation ---\\
having nothing, non-proliferating, they say.
\end{versebox}

All of these --- every word --- are synonyms for nirvana.\footnote{Pali: \textit{``Asaṅkhataṃ anataṃ anāsavañca \ldots akiñcanaṃ appapañcanti vutta''nti. Dhammānussatiyā etaṃ vevacanaṃ.''} Five stanzas, approximately forty synonyms for nirvana (\textit{nibbāna}). The Netti treats this entire list as a single entry in the Synonym Mode: one reality, forty names. Each name illuminates a different facet --- ``unborn'' emphasizes non-arising, ``peace'' emphasizes the cessation of agitation, ``island'' emphasizes safety, ``the immeasurable'' emphasizes transcendence of all quantity.}

The Community: practicing the good way, the straight way, the right way, the proper way --- that is, the four pairs of persons, the eight types of individuals. Worthy of offerings, worthy of hospitality, worthy of gifts, worthy of reverential salutation, an unsurpassed field of merit for the world. Accomplished in virtue, in concentration, in wisdom, in liberation, in the knowledge and vision of liberation. The cream of beings, the choice of beings, the extract of beings, the pillar of beings, the fragrant flower of beings.\footnote{Pali: \textit{``Suppaṭipanno ujuppaṭipanno ñāyappaṭipanno sāmīcippaṭipanno \ldots sattānaṃ sāro sattānaṃ maṇḍo sattānaṃ uddhāro sattānaṃ esikā sattānaṃ surabhipasūnaṃ.''} The standard saṅghānussati formula, extended with botanical and architectural metaphors: ``cream'' (\textit{maṇḍa}), ``extract'' (\textit{uddhāra}), ``pillar'' (\textit{esikā}), ``fragrant flower'' (\textit{surabhipasūna}).}

And virtue itself gets synonyms: an ornament (for its supreme adornment), a treasure (for overcoming all misfortune), a craft (for hitting the mark), a boundary (for not being transgressed), wealth (for ending poverty), a mirror (for seeing the teaching reflected), a tower (for the view it affords), and something that follows you everywhere --- because virtue, with its final goal of the deathless, accompanies you across every realm of existence.\footnote{Pali: \textit{``Alaṅkāro ca sīlaṃ uttamaṅgopasobhaṇatāya, nidhānañca sīlaṃ sabbadobhaggasamatikkamanaṭṭhena, sippañca sīlaṃ akkhaṇavedhitāya, velā ca sīlaṃ anatikkamanaṭṭhena, dhaññañca sīlaṃ daliddopacchedanaṭṭhena, ādāso ca sīlaṃ dhammavolokanatāya, pāsādo ca sīlaṃ volokanaṭṭhena, sabbabhūmānuparivatti ca sīlaṃ amatapariyosānan'ti.''} Eight metaphors for virtue (\textit{sīla}), each with its own rationale: (1) ornament → supreme adornment; (2) treasure → overcoming all misfortune; (3) craft → hitting the mark; (4) boundary → not transgressed; (5) wealth → ending poverty; (6) mirror → reflecting the teaching; (7) tower → providing a view; (8) universal companion → follows everywhere, ending in the deathless.}

All synonyms. One teaching, looked at from every direction. That is why the Venerable Mahākaccāyana said: ``Many are the synonyms.''\footnote{Pali: \textit{``Tenāha āyasmā mahākaccāyano `vevacanāni bahūnī'ti.''}}
